\documentclass[11pt]{article}
\usepackage[margin=1in]{geometry}
\usepackage{amsmath,amssymb,amsthm}
\usepackage[T1]{fontenc}
\usepackage{lmodern}
\usepackage[hidelinks]{hyperref}
\usepackage{xcolor}

\newcommand{\R}{\mathbb{R}}
\newcommand{\Pp}{\mathbb{P}}
\newcommand{\gammam}{\gamma}

\title{Talagrand's Convexity and Subgaussian-Vector Problems\\
       Explicitly Answered by arXiv:2605.10908}
\author{Literature-status packet for arXiv:2602.22342}
\date{22 July 2026}

\begin{document}
\maketitle

\begin{center}
\fcolorbox{black}{gray!8}{\parbox{0.88\linewidth}{
\textbf{Status: literature\_already\_answered.}
This packet records an explicit answer in a separate later paper.  It does
not claim a new proof produced by the run.}}
\end{center}

\section*{Source and exact questions}

The source is Antoine Song, \emph{Sum of Gaussian vectors and large sets},
arXiv:2602.22342v2 \cite{Song}.  Problem 0.1 on arXiv PDF page 1 asks whether
there is a universal integer $q$ such that, for every $n$ and every closed
$A\subset\R^n$ with $\gammam_n(A)\ge 2/3$, there is a convex body $K$ with
\[
  \gammam_n(K)\ge\frac12,
  \qquad K\subset \underbrace{A+\cdots+A}_{q\text{ times}}.
\]
Problem 0.2 on the same page asks whether there is a universal integer $q'>0$
such that every centered 1-subgaussian random vector $X\in\R^n$ can be
represented as
\[
  X=G_1+\cdots+G_{q'},
\]
where each $G_j$ is a standard Gaussian vector (the $G_j$ need not be
independent).  Theorem 0.3 on page 2 states that these problems are equivalent.

The current source PDF already contains an ``Added May 2026'' note on page 2
identifying a subsequent paper as a solution of both problems.

\section*{Explicit answering paper}

The separate paper Dongming (Merrick) Hua, Antoine Song, and Stefan Tudose,
\emph{On Talagrand's Convexity Conjecture}, arXiv:2605.10908v2
\cite{HST}, restates both problems, cites \cite{Song}, and says explicitly that
it resolves them affirmatively.

Its Theorem 1.1 on arXiv PDF page 2 proves the stronger structural statement:
if $X$ is dominated in convex order by a standard Gaussian vector $G$, then
there are three standard Gaussian vectors $G_1,G_2,G_3$ such that
\[
  X=G_1+G_2+G_3.
\]
Corollary 1.2 combines this with a universal convex-order comparison for
centered 1-subgaussian vectors.  The paragraph immediately following the
corollary explains how to remove the universal scaling by using a universal
finite number of Gaussian summands.  This is exactly an affirmative answer to
Problem 0.2 of \cite{Song}.  By the equivalence in Theorem 0.3 of \cite{Song},
the convexity problem is also answered affirmatively; the answering paper also
states a quantitative geometric version as Corollary 1.3.

\section*{Classification and scope}

The identification is author-explicit, not an implication inferred only by
this run: the answering paper includes the source author, cites the source,
repeats the two questions, and states that they are resolved.  It therefore
belongs in \texttt{solutions/literature\_already\_answered/}, not in an
original-solution bucket.

The full existential questions extracted from arXiv:2602.22342 are settled.
The supporting paper separately notes that a constructive proof of its
quantitative convex-body corollary remains open; that refinement is not part of
the extracted questions classified here.

\section*{Evidence files}

The packet includes the source PDF as \texttt{source\_paper.pdf}.  It includes
the answering PDF as \texttt{supporting\_paper\_2605.10908.pdf}.  The decisive
locations are source pages 1--2 and supporting pages 1--2.

\begin{thebibliography}{9}
\bibitem{Song}
A. Song,
\emph{Sum of Gaussian vectors and large sets},
arXiv:2602.22342v2, 2026.

\bibitem{HST}
D. M. Hua, A. Song, and S. Tudose,
\emph{On Talagrand's Convexity Conjecture},
arXiv:2605.10908v2, 2026.
\end{thebibliography}

\end{document}
